\section{Introduction}

Text-to-image (T2I) models have achieved impressive progress in generating high-quality images from natural language prompts~\citep{flux2024, sd35Medium, Wu2025QwenImageTR}. As these systems move into real-world applications such as marketing, product design, and slide creation, they are increasingly expected to solve practical visual tasks rather than merely render prompts.

Despite their generative ability, current T2I models remain limited on real-world tasks~\citep{He2026MindBrushIA}. A key reason is a structural mismatch between training and deployment: models are optimized for fully specified prompts~\citep{Wu2025QwenImageTR}, while real-world requests are often underspecified. In practice, successful generation may require inferring implicit user intent, retrieving up-to-date knowledge or visual references from web, and incorporating interaction history.

We refer to this mismatch as the \textbf{Context Gap}: the gap between the provided \textbf{\textit{user context}} and the \textbf{\textit{generation context}} required for T2I models. This gap motivates a paradigm shift from traditional \textbf{\textit{direct image generation}} to \textbf{\textit{agentic image generation}}, where the system must identify missing context, acquire it, and use it effectively during generation. Recent work has explored components such as plan~\citep{Yao2026PhotoAgentAP}, reason~\citep{He2026MindBrushIA}, search and tool use~\citep{Ye2026AgentBH,Feng2026GenSearcherRA,He2026MindBrushIA}, memory~\citep{He2026GEMSAM}, and self feedback~\citep{Jiang2026GenAgentST,Wang2025ImAgentAU}, but these efforts remain fragmented and do not provide a unified framework for context-centered generation.

To this end, we propose \textbf{Qwen-Image-Agent}, a unified agentic framework that integrates plan, reason, search, memory and feedback in a context-centric manner. Rather than treating user context as the final generation condition, our pipeline progressively constructs the full generation context through Context-Aware Planning and Context Grounding. Specifically, \textbf{Context-Aware Planning} operates at three levels: Information-level Planning identifies missing information and routes it to appropriate grounding strategies, Content-level Planning assembles grounded context into a detailed generation specification, and Generation-level Planning allocates context in multi-image and multi-turn scenarios. \textbf{Context Grounding} collects missing context from multiple sources, including reasoning for implicit intent inference, search for factual knowledge and visual references, memory for historical and personalized context, and feedback for iterative refinement. Overall, Qwen-Image-Agent is training-free, compatible with existing image generators, and supports both multi-image and multi-turn interaction.

Existing evaluations mainly emphasize rendering abilities~\citep{Ghosh2023GenEvalAO, Hu2024ELLAED} or isolated knowledge and reasoning abilities~\citep{Niu2025WISEAW, Zhao2025EnvisioningBT}, but fail to systematically assess the capabilities required for agentic image generation. To fill this gap, we introduce \textbf{Image-Agent-Bench (IA-Bench)}, a benchmark that evaluates four core agentic capabilities: Plan, Reason, Search, and Memory, over 17 real-world tasks, 730 test instances, and 1801 fine-grained binary checklist items. Each task is paired with a structured VLM-based evaluation protocol for reliable assessment.

Experiments on IA-Bench and prior benchmarks, including WISE-Verified~\citep{Niu2025WISEAW} and MindBench~\citep{He2026MindBrushIA}, show that Qwen-Image-Agent substantially outperforms strong agentic baselines and achieves state-of-the-art results. Ablation studies further verify the complementary benefits of different grounded contexts. Our contributions are summarized as follows:
\begin{itemize}[leftmargin=*, itemsep=2pt, topsep=2pt, parsep=0pt, partopsep=0pt]
\item We identify the \textbf{Context Gap}, i.e., the mismatch between user context and generation context as a fundamental challenge in real-world image generation. This provides a unified lens for understanding why current T2I systems fail in practical settings.
\item We propose \textbf{Qwen-Image-Agent}, a unified and context-centric framework for agentic image generation that addresses the context gap through plan, reason, search, memory and feedback.
\item We introduce \textbf{IA-Bench}, a benchmark for systematically evaluating agentic image generation along four capabilities: Plan, Reason, Search, and Memory.
\item Experiments show that Qwen-Image-Agent substantially outperforms strong agentic baselines, and achieve state-of-the-art performance on IA-Bench, Mindbench and WISE-Verified.
\end{itemize}